\section{Random Graph Theory}

\subsection{Chromatic number of Gnp}
The following is a very nice theorem by Béla!
\begin{theorem} [Béla 87]
	\label{trm:chromatic_gnhalf}
	\begin{equation}
		\chi(G(n,1/2)) = (1/2 + o(1))\frac{n}{\log_2(n)}
	\end{equation}
\end{theorem}
The lowerbound follows easily from the fact that $\alpha(G(n,1/2))
\leq 2\log_2(n)$ whp [\ref{lemma:independence_number_gnp}] and that
for any $n$-vtx graph $G$, $\chi(G) \geq n/\alpha(G)$.

The upperbound is a very nice application of Janson's
[\ref{trm:janson_ineq}]. We will just grab independent
sets of size $(2 - \eps)\log_2(n)$ while we can!.

\begin{lemma}
	For $\eps > 0$, with very high probability, every  set of vertices
	of $G(n,1/2)$ with size $m \geq n/(\log n)^2$,
	has an independent set of size $(2 - \eps)\log_2{n}$.
\end{lemma}
\begin{proof}
	We will do an union bound over $m$ sets. For that, set $k =
	(2-\eps)$ and fix a specific set $M$ of size $m$, note that
	\[
		\Pr(\alpha(G(n,1/2)[M]) < k) = \Pr(\alpha(G(m,1/2) < k)).
	\]
	Now, what is $\Pr(\alpha(G(m,1/2) < k))$? Let $S \in \binom{[m]}{k}$
	and denote $B_S = \{S \text{ is indep.}\}$.
	It is clear now that,
	\[
		\Pr(\alpha(G(m,1/2)) < k) = \Pr\bigg(\bigcap_S B_S^c\bigg)
	\]
	and we are in the hypothesis of Janson. Note that
	\[
		\mu = \sum_{S \in \binom{[m]}{k}} \Pr(S \text{ is indep.}) =
		\binom{m}{k}\bigg(\frac{1}{2}\bigg)^{\binom{k}{2}}
		\geq \bigg(\frac{m}{k}2^{-{k}/2}\bigg)^k \gg (n^{\eps})^{\log n}
	\]
	and
	\[
		\Delta = \sum_{\substack{S,T \in \binom{[m]}{k}\\S \cap T > 2}}
		\Pr(S\cup T \text{ indep.})
		= \sum_{l = 2}^{k-1}\binom{m}{k}\binom{k}{l}\binom{m - k}{k -
		l}\bigg(\frac{1}{2}\bigg)^{2\binom{k}{2} - \binom{l}{2}}.
	\]
	Now it's a bit ugly to see but this thing here is convex and the sum
	is majorized by the last two elements. We get
	\begin{align*}
		\Delta &\simeq \binom{m}{k}\binom{k}{2}\binom{m - k}{k -
		2}\bigg(\frac{1}{2}\bigg)^{2\binom{k}{2} - 1} + \binom{m}{k}k (m -
		k)\bigg(\frac{1}{2}\bigg)^{\binom{k}{2}}\\
		&\simeq 2\binom{m}{k}^2 \frac{k^4}{m^2} 2^{-2\binom{k}{2}} + km\mu
		=  \frac{2k^4}{m^2}\mu^2 + km\mu \leq \frac{3k^4}{n^2}\mu^2.
	\end{align*}
	By the second part of Janson, we get that
	\[
		\Pr\bigg(\bigcap_S B_S^c\bigg) \leq \exp(- cn^2/k^4) \lll 2^{-n}
	\]
	So we can easily do a union bound, over all sets.
\end{proof}

\begin{proof}
	[Proof of Theorem \ref{trm:chromatic_gnhalf}]
	Just use the previous lemma and remove independent sets untill you
	have less than $n/(\log n)^2$ vertices left. At that point, do a
	trivial coloring.
\end{proof}

\subsubsection{Concentration of the Chromatic number}
We will look into a very weird type of result. We will show that with
high probability, if $p$ is in the right range,
$\chi(G(n,p))$ is super concentrated. Our main objective is to prove
the following surprising statement
\begin{theorem}[Alon-Krivilevich 95]
	\label{trm:chromatic_concentration_alon_krivilevich}
	Let $\delta > 0$, $p = n^{-1/2 - \delta}$ and $G \sim G(n,p)$ then
	there exists $t = t(n)$ s.t w.h.p
	\[
		\chi(G) \in \{t, t + 1\}.	
	\]
\end{theorem}
\begin{remark}
	This is a beautiful theorem, although an extremely weird one, since
	we don't even know the value of $t$.
\end{remark}
From now on, we will always think of $G \sim G(n,p)$, to diminish
notation clustering.
The main tool for this kind of result is Luczak's lemma. It showed up
in the first proofs of statements of this kind.

\begin{lemma}[Luczak]
	\label{lemma:luczak_martingale}
	For any $\eps > 0$, let $t = t(n,p,\eps)$ be the least integer s.t
	\[
		\Pr(\chi(G) \leq t) \geq \eps.
	\]
	Let $X$ be a r.v that counts the minimal number of vertices to be
	removed from $G$ to make it $t$-colorable. Then, if $\lambda$ is such that
	$\eps = e^{-\lambda^2/2}$, we have
	\[
		\Pr(X > 2\lambda \sqrt{n}) \leq \eps.
	\]
\end{lemma}
\begin{proof}
	We will do a clever martingale argument (note crucially here that we
	don't really know that $t$ is). $X$ satisfies the Lipchitz condition
	for the vertex exposure martingale, so that it is typically
	concentrated [\ref{trm:azuma_inequality}]. More precisely,
	\begin{align}
		\Pr(X \leq \Ex X - \lambda \sqrt{n} ) < e^{-\lambda^2/2} \quad
		\text{and} \quad \Pr(X \geq \Ex X + \lambda \sqrt{n}) < e^{-\lambda^2/2}.
	\end{align}
	In particular, by the definition of $X$, and $t$, we know that with
	probability at least $\eps$, $X = 0$. By the first equation,
	this means that $\Ex X < \lambda\sqrt{n}$. Plugging that into the
	second equation we get the result.
\end{proof}

Let's get ready for theorem
[\ref{trm:chromatic_concentration_alon_krivilevich}]. We will need
some properties of $G(n,p)$ at this range.
\begin{lemma}
	\label{chromatic:lemma:stuff_gnp_satisfies}
	[Properties of $G(n,p)$] Let $\delta > 0$ and $G \sim G(n,p)$ with $p
	= n^{-1/2 - \delta}$. The following properties all hold with high
	probability (less than some $\eps$ for big $n$).
	\begin{enumerate}[label=(\roman*)]
		\item $\Delta(G) \leq 3d = 3np$.
		\item $\chi(G) \geq d/2 \log n$.
		\item For some fixed $r = r(\delta)$  we have that for every $U
			\subset V(G)$ with $|U| \leq O(\sqrt{n})$, $e(G[U]) \leq r|U|$.
			Note that, by degeneracy,
			this makes every set $W \subset V(G)$ with $|W| \leq O(\sqrt{n})$
			$2r$-colorable.
		\item If $\delta > 1/5$, then for every set $U \subset V(G)$ with
			$|U| \leq O(n^{1-\delta})$, $e(G[U]) \leq |U|n^{1/10}$.
	\end{enumerate}
	Furthermore, let $X_{u,v}$ be the number of paths of length 3 (edges)
	between $u$ and $v$. We also find that
	\begin{enumerate}[label=(\roman*)]
			\setcounter{enumi}{4}
		\item If $\delta \geq 1/6$, we find that $\max_{u,v \in V(G)} X_{u,
			v} \leq \log n$.
		\item For $0 < \delta < 1/6$, this maximum is at most  $d^3\log n / n$.
	\end{enumerate}
\end{lemma}
\begin{proof}
	All of these follow from the first moment. (i) is just Chernoff, (ii)
	follows directly from [\ref{lemma:independence_number_gnp}]. (iii) and (iv)
	are as well simple first moment bounds. The only interesting ones are
	(v) and (vi). To show them we note that, fixed $u,v \in V(G)$,  every
	edge participates
	in at most $c_0 = 3/\delta$ paths of length 3 between $u$ and $v$  w.
	h.p. In particular, each path of length 3 between $u,v$ touches at
	most $3c_0$ others.
	So there must be $X_{u,v}/(3c_0 + 1)$ disjoint paths between $u$ and
	$v$. Using this bound, we show the result.
\end{proof}

Now we are ready to prove the theorem. It suffices to show the
following proposition.

\begin{prop}
	\label{prop:main_prop_chr_concentration}
	Suppose $G(\delta)$ is a graph on $n$ vertices satisfying the
	conclusions of the previous lemma [\ref{chromatic:lemma:stuff_gnp_satisfies}].
	Moreover, suppose, following lemma [\ref{lemma:luczak_martingale}],
	that $G \geq t \geq d/2 \log n$ and that $G$ has a set $U_0 \subset V(G)$
	s.t $|U| \leq 2\lambda\sqrt{n} = O(\sqrt{n})$ and $\chi(G[V(G)
	\setminus U]) = t$. Then $G$ is $t + 1$ colorable.
\end{prop}

\begin{remark}
	We will prove a totally deterministic statement, using probabilistic
	methods, to prove a probabilistic result! The irony.
\end{remark}

I will lay out the proof plan. The idea is to find a set $U$ with
$U_0 \subset U$, $|U| =
O(\sqrt{n})$ such that every vertex in $V(G) \setminus U$ has few
neighbors in $U$ (say $10r$ which is constant).
We know that $V(G) - U$ is $t$ colorable and that $U$, by its size,
is $2r$ colorable as well. Fix any $t$-coloring of $V(G) \setminus
U$. Say $U = \{u_1, \dots, u_k\}$, set
$N(u_i) = N_G(u_i) \setminus U$, also set $N(U) = \bigcup_{j = 1}^k N(u_j)$.
We will try to exchange $2r$ colors from each neighborhood $N(u_i)$
with a new color $\{t+1\}$ as to leave $2r$ possible
color choices for the vertex $u_i$, after that we can just use
degeneracy to greedily color $U$ with $2r$ colors and finish the proof (Having
properly (t+1)-colored our graph).

Although the strategy is the same, how we choose the colors will be
different depending on $\delta$.
So for now I will only show the $\delta < 1/5$ case.

\begin{proof}[Proof of \ref{prop:main_prop_chr_concentration}]
	(For $\delta \leq 1/5$)
	Let's formalize the previous paragraph. Fix $r$ constant as in
	[\ref{chromatic:lemma:stuff_gnp_satisfies}]. To find $U$, start with
	$U = U_0$ of size $2\lambda\sqrt{n}$ and while there are vertices
	$v$ s.t. $d_U(v) > 10r$, add them to $U$. We cannot have that $|U|
	\geq 2 |U_0|$ as that would break the hypothesis that $e(G[U]) \leq r|U|$.
	So we really have $|U| = O(\sqrt{n})$ and every vertex outside of
	$U$ sends in at most constant edges.

	Now let $f : V(G) \setminus U \to \{1,\dots, t\}$ be a proper
	coloring of the rest of the world. To help us recolor the
	neighborhood of $U$, we will do a
	unneeded but very helpful step of defining an auxiliary graph $H$ of
	$N(U)$. For each $1 \leq i \leq k$, let $W_i$ be a copy of $N(u_i)$,
	we define $H$
	with vertex set $\bigsqcup_{i = 1}^k W_i$. Note that each vertex $a
	\in W$ corresponds to a vertex $a'$ in $N(U)$, we connect $a$ to $b$
	if $a' \sim b'$.
	Note that $f$ induces a proper coloring $f'$ of $H$ as well (in some
		sense we are just blowing up the vertices of $N(U)$ according to how
	many neighborhoods they belong to).
	Also, a set $A$ in $H$ is independent if and only if the
	corresponding set is independent in $N(U)$. Very importantly, note
	that each edge of $U$ corresponds to
	at most $(10r)^2$ edges in $H$.

	Having defined this graph we can be more precise about our plan of
	attack. Each $W_i$ is colored with $t$ colors.
	For each $1 \leq i \leq k$ we want to choose $2r$ colors from $\{1,
	\dots, t\}$, making a set $C_i$, such that $W_0 = \bigcup_{j = 1}^k
	\{w \in W_j : f'(w) \in C_j\}$ is independent.
	If we do so, we can just paint $W_0$ with the color $t+1$, paint the
	corresponding vertices in $N(U)$ with this color, and leave $2r$
	colors available for each $u_i$.

	We will choose this $C_i$ randomly
	(but differently depending on $\delta$). If $\delta \leq 1/5$,
	we will use the LLL [\ref{trm:lll}], if $\delta > 1/5$ we will do a
	clever counting argument. As we are in the $\delta \leq 1/5$ case,
	let's proceed.

	For each $i$ choose $C_i \subset
	\binom{[t]}{2r}$ uniformly and randomly. For each edge $e = \{w_1,
	w_2\}$ in $H$, with $w_1 \in W_{i_1}$ and $w_2 \in W_{i_2}$,
	define the bad event $B_e$ to be $\{f'(w_1) \in C_{i_1} \text{ and }
	f'(w_2) \in C_{i_2}\}$. Note that for any edge $e$,
	\[
		\Pr(B_e) \leq \bigg(\frac{2r}{t}\bigg)^2 = O((\log n/ d)^2).
	\]
	What about our dependency graph? We note that $B_{w_1,w_2}$ is
	completely independent of $B_{w_3, w_4}$ if $\{W_{i_1}, W_{i_2}\}
	\cap \{W_{i_3}, W_{i_4}\} = \varnothing$.
	That is, by revealing $C_j$ for $j \not \in \{i_1, i_2\}$ we gain no
	information about $B_{\{w_1,w_2\}}$. So what is the degree of $B_e$?
	It is the number of
	edges that touch $W_{i_1}$ or $W_{i_2}$. Because $|W_{i_1}| = N(u_i)
	\leq 3d = o(\sqrt{n})$, $e(W_i) = O(d)$. What about edges between
	$W_{i_1}$ and $W_j$? Because of the assumption over $X(u_{i_1},
	u_j)$ (number of paths of length 3 between both vtxs), and that each
	edge of $G$ is corresponds to at most $(10r)^2$ edges in $H$, we note that
	\[
		e(W_{i_1}, W_j) \leq
		\begin{cases}
			(10r)^2 \log n & \text{ if } \delta > 1/6\\
			(10r)^2 d^3\log n/n & \text{ if } \delta \leq 1/6
		\end{cases}
	\]
	Overall, this means that
	\[
		d(B_e) \leq
		\begin{cases}
			O(d) + k \cdot O(\log n) = O(n^{1/2}\log n) & \text{ if } \delta > 1/6\\
			O(d) + k \cdot O(d^3 \log n / n) = O(d^3\log n / n^{1/2}) & \text{
			if } \delta \leq 1/6\\
		\end{cases}
	\]
	In both cases, as $\delta \leq 1/5$, $d \geq n^{2/5}$, we get that
	\[
		\Pr(B_e) \cdot d(B_e) = o(1)
	\]
	and by the LLL [\ref{trm:lll}] we get our result.
\end{proof}

Now we delve into what should be the easier case (closer to what
people already had solved) of $\delta > 1/5$.
In our computations we will need to use that the result was already
proven for $\delta > 3/8$, which is a one page proof
by Luczak [\url{https://link.springer.com/article/10.
1007/BF01205080}]. It is worth reading it, it has some similar ideas.

Our main goal will be the same, still choose for
each $W_i$ a set of colors $C_i$, such that the set $W_0$ as defined
before is an independent
set. But here we will have to do some alteration, we will take each
$C_i$ in a way that $C_j$ and $C_i$ are compatible for $j < i$,
furthermore, we will take each $C_i$ a bit bigger so that we can
remove the colors of the edges inside $W_i$, which are few.
Let's engage the proof now.

\begin{proof}[Proof of \ref{prop:main_prop_chr_concentration}]
	(For $1/5 < \delta \leq 3/8$) As before define $U$, $H$ and its
	induced coloring $f'$ with $t \geq d/2\log n$ colors.
	By Lemma \ref{chromatic:lemma:stuff_gnp_satisfies}, for each $s \leq
	k$, the sets $N(u_{i_1}), \dots, N(u_{i_s})$ have size less than
	$s3d \leq n^{1 - \delta}$,
	and so span less than $3sdn^{1/10}$ edges. In particular, we can
	always order the $u_i$'s in such way that that $N(u_i)$ sends at most
	$3dn^{1/10}$ edges for all $i'<i$. It follows that, in this order,
	\[
		e\bigg(W_i, \bigcup_{j < i} W_j\bigg) \leq 3(10r)^2dn^{1/10} = O(dn^{1/10})
	\]
	as well. For each $i$, we will choose a random $2r + \log n$ set of
	colors from the ones that are "available".
	A color $c \in \{1,\dots,t\}$ is not available if there exists an $j
	< i$ and an edge $(w_j, w_i)$ between $W_j, W_i$ with $f'(w_j) = c$.
	Call $x_i$ the probability that for some $i'\leq i$, less than $t/2$
	colors are available when choosing $C_{i'}$, we will try to show that
	$x_k < 1$. Obviously, $x_1 = 0$, now we note that, because $\delta
	\geq 1/6$, $e(W_j, W_i) \leq O(\log n)$, so each color
	chosen in $C_j$ may kill at most $\log n$ colors when choosing $C_i$.
 Because we need to kill $t/2$ colors, we get that
	\begin{align*}
		x_i &= x_{i-1} + (1 - x_{i-1})\binom{O(dn^{1/10})}{\frac{t}{2O(\log
		n)}} \bigg(\frac{2r + \log n}{t/2}\bigg)^{\frac{t}{2O(\log n)}}\\
		&\leq x_{i-1} + e^{-n^{1/10}}
	\end{align*}
	So clearly, $x_k \leq \sqrt{n}e^{-n^{1/10}} \to 0$ (in this
	calculation it seems we use $\delta \leq 3/8$ somewhere).
	Now, this means we can choose the colors $C_i$ such that they are
	compatible. Inside each $W_i$ there are at most $\log n$ edges,
	kill the colors corresponding to them, we are still left with $2r$
	colors for each $C_i$ that make $W_0$ an independent set.
	The proof follows the same way as before now.
\end{proof}
